% p316_1; source PDF page 316; manually verified editable table.
\begin{tikzpicture}[x=1cm,y=-1cm,inner sep=0pt,outer sep=0pt]
\path[use as bounding box] (0,0) rectangle (16.000,10.220);
\draw[line width=.4pt] (0.000,0.000) rectangle (2.700,0.805);
\node[align=center,text width=2.440cm,font={\linespread{1}\fontsize{10}{13.50}\selectfont\bfseries}] at (1.350,0.403) {危害公共安全};
\draw[line width=.4pt] (2.700,0.000) rectangle (6.400,0.805);
\node[align=center,text width=3.440cm,font={\linespread{1}\fontsize{10}{13.50}\selectfont\bfseries}] at (4.550,0.403) {要点};
\draw[line width=.4pt] (6.400,0.000) rectangle (16.000,0.805);
\node[align=center,text width=9.340cm,font={\linespread{1}\fontsize{10}{13.50}\selectfont\bfseries}] at (11.200,0.403) {范例};
\draw[line width=.4pt] (0.000,0.805) rectangle (2.700,4.210);
\node[align=center,text width=2.440cm,font={\linespread{1}\fontsize{10}{13.50}\selectfont}] at (1.350,2.507) {安全};
\draw[line width=.4pt] (2.700,0.805) rectangle (6.400,4.210);
\node[align=left,text width=3.440cm,font={\linespread{1}\fontsize{10}{13.50}\selectfont}] at (4.550,2.507) {安全内容是指人身伤亡、财产损失（物质性损害），不包括精神损害、心理恐慌（非物质性损害）};
\draw[line width=.4pt] (6.400,0.805) rectangle (16.000,4.210);
\node[align=left,text width=9.340cm,font={\linespread{1}\fontsize{10}{13.50}\selectfont}] at (11.200,2.507) {甲将面粉邮寄给某单位，谎称是炭疽热病菌粉，引起众人恐慌，乱作一团。甲构成投放虚假危险物质罪。该罪不是危害公共安全犯罪，而是妨害社会管理秩序犯罪，因为引起心理恐慌，不属于危害公共安全};
\draw[line width=.4pt] (0.000,4.210) rectangle (2.700,10.220);
\node[align=center,text width=2.440cm,font={\linespread{1}\fontsize{10}{13.50}\selectfont}] at (1.350,7.215) {公共};
\draw[line width=.4pt] (2.700,4.210) rectangle (6.400,10.220);
\node[align=left,text width=3.440cm,font={\linespread{1}\fontsize{10}{13.50}\selectfont}] at (4.550,7.215) {公共是指公众（多数人或不特定人）。危害公众安全，是指\CJKunderwave{危险具有危及多数人的可能性}};
\draw[line width=.4pt] (6.400,4.210) rectangle (16.000,6.815);
\node[align=left,text width=9.340cm,font={\linespread{1}\fontsize{10}{13.50}\selectfont}] at (11.200,5.512) {第一，虽然死亡人数有三人，但是危险的危及范围非常确定，只会危及特定少数人，不属于危害公共安全。例如（灭门案），甲向乙家锅里投毒，杀死了乙一家三口，不构成投放危险物质罪，只构成故意杀人罪};
\draw[line width=.4pt] (6.400,6.815) rectangle (16.000,10.220);
\node[align=left,text width=9.340cm,font={\linespread{1}\fontsize{10}{13.50}\selectfont}] at (11.200,8.517) {第二，虽然死亡人数只有一人，但是危险具有危及多数人的可能性，属于危害公共安全。例如（食堂锅里投毒案），甲是校园食堂师傅，向一锅汤里放老鼠药，一名学生中毒死亡。由于危险的危及范围可大可小，所以甲制造了公共危险，危害了公共安全，构成投放危险物质罪①};
\end{tikzpicture}
